\section{坏小孩、刺猬和新物种：硅谷资本的创始人 taste}

上一章讲 Robinhood 的业务，本章讲创始人和资本 taste。节目里有几个生动词：乖小孩、坏小孩、刺猬型、反叛者、哪吒型。它们不是为了贴标签，而是描述硅谷资本如何识别“能打破常规”的 founder-led 公司。

\begin{figure}[H]
\centering
\includegraphics[width=0.96\textwidth]{founder-archetypes.png}
\caption{硅谷创始人画像：乖小孩、坏小孩、刺猬、反叛者、哪吒与 Founder-led。自制概念图，依据 00:40:29--00:44:26 与 00:57:13--00:58:22 对谈内容整理。}
\end{figure}

\begin{knowledgebox}{读图：这些不是性格测试，而是投资 taste}
乖小孩代表合规稳健，坏小孩代表敢打破规则，刺猬型代表单点极深，反叛者代表挑战旧秩序，哪吒型代表不按牌理出牌。资本真正寻找的是 founder-led 的持续执行力，而不是为了叛逆而叛逆。
\end{knowledgebox}

\subsection{颠覆性公司的共性}

本节从创始人标签回到公司基本面。一个创始人“坏”或“反叛”本身没有价值，只有当这种性格能服务于大市场、强产品和快速执行时，才会成为资本愿意押注的特质。

Freda 总结颠覆性公司有几个共性。第一，TAM 足够大，例如 Robinhood 想做整个金融行业超级应用。第二，产品足够好，最好不需要大量销售也能高速增长。第三，天时地利人和：智能手机和带宽成熟催生短视频，摩尔定律终结和 GPU 推动 AI，下沉市场空白催生拼多多。很多颠覆者在刚出现时都不合常理：Robinhood 零佣金，OpenAI 疯狂烧钱，Netflix 改变 DVD 租赁，TikTok 重写内容分发。

\begin{importantbox}{颠覆性不是“奇怪”本身}
一个公司看起来不合常理，只是必要不充分条件。它还必须面对巨大市场、拥有极强产品、踩中技术和社会时机，并能用组织执行力把反常识变成新常识。
\end{importantbox}

\subsection{预测市场：一个新物种}

接下来用预测市场作为“新物种”样本。它既是金融产品，也是信息发现机制，还和散户参与、监管宽松、媒体信任危机共同作用。

节目中提到的 prediction market（预测市场）是“新物种”例子。Kalshi、Polymarket 等平台让用户用小额资金下注事件：大选结果、财报超预期、模型发布时间，甚至娱乐事件。Freda 认为这不是小众市场，因为交易额增长很快，Bloomberg 和 Google Finance 都开始接入数据。

预测市场出现的背景包括：美国主流媒体信任危机、疫情/关税/政治动荡带来的不确定性、监管相对宽松、以及人们相信“总有人知道一点什么”。它和 Robinhood 的关系在于，预测市场既是金融产品，也是大众参与事件定价的新入口。

\subsection{本章小结}

硅谷资本喜欢的不是“听话公司”，而是能在大市场里用强产品、强执行和反常识判断形成新物种的公司。创始人画像只是外壳，真正重要的是市场、产品、时机和组织能力。

